\chapter{Podstawowe pojęcia}

Głównym celem pracy jest omówienie form normalnych macierzy zespolonych symetrycznych, antysymetrycznych oraz ortogonalnych.
Dlatego najpierw omówimy odpowiednie przypadki rzeczywiste macierzy symetrycznych, antysymetrycznych oraz ortogonalnych.
Formą normalną będziemy nazywać postać diagonalną macierzy, a w przypadkach gdy nie jest ona osiągalna, najprostszą postać blokowo-diagonalną, do jakiej macierz można sprowadzić. 

W tym rozdziale będą przedstawione podstawowe pojęcia, a w dalszej części formy normalne dla macierzy rzeczywistych symetrycznych, antysymetrycznych oraz ortogonalnych.

\section{Przestrzeń unitarna}

Niech $V$ będzie przestrzenią liniową nad ciałem $\K$ liczb zespolonych lub rzeczywistych ($\K$ = $\C$ lub $\K$ = $\R$). Przekształcenie $\inner{\cdot}{\cdot} \colon V \times V \to \K$ przyporządkowujące parze wektorów ($x$, $y$) $\in$ $V \times V$ liczbę z $\K$ nazywamy \textbf{iloczynem skalarnym}, jeżeli dla dowolnych wektorów $x$, $y$, $z$ $\in$ $V$ spełnione są następujące warunki:

\begin{enumerate}[leftmargin=1cm, itemsep=2pt, label=(\arabic*)]
    \item $\inner{x}{y}$ = $\overline{\inner{y}{x}}$, gdzie $\overline{\inner{y}{x}}$ jest sprzężeniem zespolonym liczby $\inner{y}{x}$
    \item $\inner{ax}{y}$ = $a\inner{x}{y}$, gdzie $a \in \K$
    \item $\inner{x + y}{z}$ = $\inner{x}{z} + \inner{y}{z}$
    \item $\inner{x}{x} \geq 0$ oraz $\inner{x}{x} = 0 \iff x = 0$ (warunek o dodatniej określoności)
\end{enumerate}
Z (1), (2), (3) wynika, że:
\begin{enumerate}[leftmargin=1cm, itemsep=2pt, label=(\arabic*), start=5]
    \item $\inner{x}{ay}$ = $\overline{\inner{ay}{x}}$ = $\overline{a}\inner{x}{y}$, gdzie $a \in \K$
    \item $\inner{x}{y + z}$ = $\inner{x}{z} + \inner{x}{y}$
\end{enumerate}

\begin{definition}
Przestrzeń liniową $V$ nad ciałem liczb $\K$ wraz z wyróżnionym iloczynem skalarnym $\inner{\cdot}{\cdot}$, tj. parę ($V$, $\inner{\cdot}{\cdot}$) nazywamy $\textbf{przestrzenią unitarną}$.
W przypadku gdy $\K = \R$, przestrzeń unitarną nazywamy \textbf{przestrzenią euklidesową}.
\end{definition}

Dalej będziemy rozpatrywać tylko przestrzenie o skończonym wymiarze.

\begin{definition}
    \textbf{Normą wektora} $x \in V$ nazywamy liczbę $\norm{x} = \sqrt{\inner{x}{x}}$ 
\end{definition}

Niech $V$ będzie $n$-wymiarową przestrzenią unitarną z dowolną bazą $\base{e}$. Niech \textbf{x}, \textbf{y} będą dowolnymi wektorami z $V$, gdzie $x_i$, $y_i$ są współrzędnymi tych wektorów:

\begin{center}
    $$\textbf{x} = \sum_{i=1}^{n}{x_{i}e_{i}} \quad \textbf{y} = \sum_{i=1}^{n}{y_{i}e_{i}}$$
\end{center}
Wtedy z (2), (3), (5), (6):
\begin{center}
    $$\inner{\textbf{x}}{\textbf{y}} = \inner{
        \sum_{i=1}^{n}{x_{i} \textbf{e}_i}
    }{
        \sum_{k=1}^{n}{y_{k} \textbf{e}_k}
    } = \sum_{i=1}^{n}{\sum_{k=1}^{n}{x_{i} \bar{y}_{k} \inner{\textbf{e}_{i}}{\textbf{e}_{k}}}}$$
\end{center}
Niech
\begin{equation}\label{eq:def_h}
    h_{ik} = \inner{\textbf{e}_i}{\textbf{e}_k}, \text{ dla }k,i = 1, \ldots, n 
\end{equation}
Wtedy: 
\begin{equation}\label{eq:std_inner}
    \inner{\textbf{x}}{\textbf{y}} = \sum_{i,k=1}^{n}{h_{ik} x_{i} \bar{y}_{k}}
\end{equation}
W szczególności:

\begin{equation}
    \inner{\textbf{x}}{\textbf{x}} = \sum_{i,k=1}^{n}{h_{ik} x_{i} \bar{x}_{k}}
\end{equation}
gdzie z \eqref{eq:def_h} oraz własności (1), wynika, że $h_{ik} = \bar{h}_{ki}$.

W ten sposób zdefiniowaliśmy iloczyn skalarny dla bazy $\textbf{e}_1, \ldots, \textbf{e}_n$. 

\begin{definition}
    Układ wektorów ($x_1, \ldots, x_n$) który spełnia:
    $$ a_1x_1 + a_2x_2 + \ldots + a_nx_n = 0 \iff a_1 = a_2 = \ldots = a_n = 0$$
    będziemy nazywać liniowo niezależnym, tj. żaden z wektorów nie jest kombinacją liniową wektorów pozostałych.
\end{definition}
\begin{definition}
    Dwa wektory $x, y \in V$ będziemy nazywać ortogonalnymi lub prostopadłymi jeśli $\inner{x}{y} = 0$, piszemy wówczas $x \perp y$.
\end{definition}
\begin{definition}
    Układ wektorów ($x_1, \ldots, x_n$) w przestrzeni unitarnej $V$ będziemy nazywać ortogonalnym jeśli $x_i \perp x_j$ dla $i \ne j$, gdzie $i,j = (1, \ldots, n)$.
\end{definition}
\begin{theorem}
    Jeśli układ wektorów ($x_1, \ldots, x_n$) jest ortogonalny, to jest liniowo niezależny.
\end{theorem}
\begin{proof}
    Załóżmy że:
    $$ a_1x_1 + \ldots + a_nx_n = 0$$
    Wymnożomy obustronie przez dowolny $x_k$, gdzie $k = (1, \ldots, n)$:
    $$ \inner{a_1x_1 + \ldots + a_nx_n}{x_k} = \inner{a_1x_1}{x_k} + \ldots + \inner{a_nx_n}{x_k} = \sum_{i=1}^{n}{\inner{a_ix_i}{x_k}} = 0 $$
    Ponieważ każdy wektor wśród ($x_1, \ldots, x_n$) jest ortogonalny:
    $$ \sum_{i=1}^{n}{\inner{a_ix_i}{x_k}} = \sum_{i=1}^{n}{a_i\inner{x_i}{x_k}} = \inner{a_kx_k}{x_k} = a_k\inner{x_k}{x_k} = a_k\norm{x_k}^2 = 0 $$
    Ponieważ $\norm{x_k}^2 > 0$, więc $a_k = 0$.
\end{proof}

Jeżeli baza $\textbf{e}_1, \ldots, \textbf{e}_n$ jest ortonormalna, to znaczy taka, że każdy wektor bazy jest ortogonalny do pozostałych oraz ma normę jeden:
\begin{equation}\label{eq:def_ortonormal_base}
    \inner{\textbf{e}_i}{\textbf{e}_k} = \delta_{ik} = \begin{cases}
    0, \text{ dla } i \neq k \\
    1, \text{ dla } i = k
\end{cases}
\end{equation}
Wtedy dla bazy ortonormalnej iloczyn skalarny przyjmuje postać, którą będziemy nazywać \textbf{standardową}:
\begin{equation}\label{eq:std_normal_inner}
    \inner{\textbf{x}}{\textbf{y}} = \sum_{i=1}^{n}{x_{i} \bar{y}_{i}}
\end{equation}
\begin{equation}
    \inner{\textbf{x}}{\textbf{x}} = \sum_{i=1}^{n}{x_{i} \bar{x}_{i}} = \sum_{i=1}^{n}{|x_{i}|^2}
\end{equation}
Dla wektorów \textbf{x}, \textbf{y} ze współrzędnymi $x_i$ i $y_i$ dla $i = 1, \ldots, n$ w bazie $\textbf{e}_1, \ldots, \textbf{e}_n$
przestrzeni euklidesowej iloczyn skalarny redukuje się do:
\begin{equation}\label{eq:std_euclid_inner}
    \inner{\textbf{x}}{\textbf{y}} = \sum_{k,i=1}^{n}{h_{ik} x_{i} {y}_{k}}
\end{equation}
Dla bazy ortonormalnej:
\begin{equation}\label{eq:std_euclid_normal_inner}
    \inner{\textbf{x}}{\textbf{y}} = \sum_{i=1}^{n}{x_{i} {y}_{i}}
\end{equation}
\begin{equation}
    \inner{\textbf{x}}{\textbf{x}} = \sum_{i=1}^{n}{x_{i}^2}
\end{equation}

\

\begin{definition}
    Sprzężeniem hermitowskim macierzy $A$ nazywamy macierz powstałą przez złożenie transpozycji oraz sprzężenia zespolonego elementów, i oznaczamy jako:
    $$ A^* = \overline{A^T} $$ 
\end{definition}

\begin{definition}
    Macierzą hermitowską będziemy nazywać taką macierz $H$, która spełnia warunek:
    $$H = H^*$$
\end{definition}

\begin{definition}
    Macierzą symetryczną będziemy nazywać taką macierz $S$, która spełnia warunek:
    $$S = S^T$$
\end{definition}

\begin{definition}
    Macierzą antysymetryczną będziemy nazywać taką macierz $K$, która spełnia warunek:
    $$K = -K^T$$
\end{definition}

\begin{definition}
    Macierzą unitarną będziemy nazywać taką macierz $U$, która spełnia warunek:
    $$U^*U = UU^* = I,$$
    gdzie $I$ macierz identycznościowa.
\end{definition}

\begin{definition}
    Macierzą ortogonalną będziemy nazywać taką macierz $O$, która spełnia warunek:
    $$O^TO = OO^T = I,$$
    gdzie $I$ macierz identycznościowa.
\end{definition}

\begin{definition}
Odwzorowanie $A \colon X \to Y $, gdzie $X$ i $Y$ są przestrzeniami liniowymi nad ciałem liczb $\K$, nazywa się \textbf{operatorem liniowym} jeżeli spełnia następujące warunki dla dowolnych dwóch wektorów $x, y \in X$ oraz dowolnego skalara $a \in \K$:
\begin{enumerate}[leftmargin=1cm, itemsep=2pt, label=(\arabic*)]
    \item $A(x + y) = A(x) + A(y)$ (zachowuje dodawanie wektorów)
    \item $A(ax) = aA(x)$ (zachowuje mnożenie przez skalar)
\end{enumerate}
\end{definition}

\begin{definition}[Wartość i wektor własny]
    Niech $V$ będzie przestrzenią liniową nad ciałem liczb $\K$, oraz $A \colon V \to V$ jest operatorem liniowym, który odwzorowuje przestrzeń $V$ samą w siebie. Jeśli istnieje taki niezerowy wektor $v \in V$ dla którego jest spełniony warunek:
    $$ Av = \lambda v,$$
    gdzie $\lambda \in \K$ pewien skalar, to $v$ nazywa się \textit{wektorem własnym} oraz $\lambda$ nazywa się \textit{wartością własną} przekształcenia $A$.
\end{definition}

Jeśli przestrzenie liniowe $X$ i $Y$ mają skończony wymiar $\dim X = n$, $\dim Y = m$, oraz ustalimy w $X$ bazę $\mathbf{e}_1, \ldots, \mathbf{e}_n$, a w $Y$ bazę $\mathbf{g}_1, \ldots, \mathbf{g}_m$, to odwzorowaniu liniowemu $A$ odpowiada macierz $m \times n$ w następującym sensie: jeżeli wektorowi $\mathbf{x} = \sum_{i=1}^{n} x_i \mathbf{e}_i \in X$ operator $A$ przyporządkowuje wektor $\mathbf{y} = \sum_{i=1}^{m} y_i \mathbf{g}_i \in Y$, to współrzędne tych wektorów są związane wzorem $\mathbf{y} = A\mathbf{x}.$
\begin{remark}
W dalszej części pracy, ustaliwszy bazę w przestrzeni $X$ i $Y$, będziemy utożsamiać operator liniowy $A \colon X \to Y$ z jego macierzą $A$ w tej bazie, pisząc po prostu $A$ zarówno dla operatora, jak i dla odpowiadającej mu macierzy.
\end{remark}